% Research review (English). Build: make research → thesis/build/research-en.pdf
% Synthesis of reports/, research_notes/, docs/pedagogy/, docs/research/. Sources — SRC-NNN.
\documentclass[a4paper,11pt]{article}
\newcommand{\mainlang}{english}
\newcommand{\docstyle}{article}
% Преамбула: шрифты (латиница, кириллица, арабский с огласовками), поля, библиография.
\usepackage{etoolbox}
\usepackage{fontspec}
\usepackage{polyglossia}
\setdefaultlanguage{\mainlang}
\ifdefstring{\mainlang}{english}{\setotherlanguages{russian,arabic}}{\setotherlanguages{english,arabic}}

% Times New Roman есть в macOS; на другой машине — Liberation Serif или Tempora.
\setmainfont{Times New Roman}
\newfontfamily\cyrillicfont{Times New Roman}
% Арабский — тот же шрифт, что рисует карточки бота (ADR-0024): огласовки и шадда не теряются.
\newfontfamily\arabicfont[Script=Arabic,Path=../src/muallim/presentation/telegram/assets/fonts/,
  Extension=.ttf,UprightFont=*-Regular,BoldFont=*-Bold,Scale=1.25]{ScheherazadeNew}
% \ar{…} — арабская вставка в строке; источник строки — в тексте рядом ([SRC-NNN]).
\newcommand{\ar}[1]{\textarabic{#1}}

\usepackage[left=3cm,right=1.5cm,top=2cm,bottom=2cm]{geometry} % [утвердить] по требованиям вуза
\usepackage{setspace}
\onehalfspacing
\usepackage{indentfirst}

\usepackage{graphicx}
\graphicspath{{../docs/architecture/diagrams/}{../docs/design/screens/}{figures/}}
\usepackage{booktabs,tabularx,longtable}
\usepackage{caption}
\usepackage{csquotes}
\usepackage[hidelinks]{hyperref}

% Библиография — тот же sources.bib, что у всего проекта (ключи SRC-NNN).
\ifdefstring{\mainlang}{english}
  {\usepackage[backend=biber,style=ieee,sorting=none]{biblatex}}
  {\usepackage[backend=biber,style=gost-numeric,sorting=none]{biblatex}}
\addbibresource{../docs/research/sources.bib}

% ID рисунков и таблиц из docs/thesis/figures.md: \label{fig:R-01}, \label{tab:T-01}.

% Данные для титульных листов всех документов. [утвердить] — с руководителем и учебным офисом.
\newcommand{\thesisProgramCode}{09.04.01}                     % [утвердить] код направления
\newcommand{\thesisProgramRu}{Информатика и вычислительная техника} % [утвердить]
\newcommand{\thesisProgramEn}{Computer Science}                % [утвердить]
\newcommand{\thesisSpecialtyRu}{Программная инженерия}
\newcommand{\thesisSpecialtyEn}{Software Engineering}
% Тема — предложение автора (2026-10-18; с ADR-0051 — с адаптацией модели), утверждает руководитель [утвердить].
\newcommand{\thesisTopicRu}{ИИ-тренажёр чтения по букварю «Муаллим ас-Сани»: адаптация открытой модели распознавания фонем к речи русскоязычных начинающих и оценка обратной связи}
\newcommand{\thesisTopicEn}{An AI Reading Tutor for the \emph{Muallim al-Sani} Primer: Adapting an Open Phoneme Recognition Model to Russian-Speaking Beginners and Evaluating Pronunciation Feedback}
\newcommand{\thesisShortRu}{ИИ-тренажёр «Муаллим ас-Сани» и адаптация модели}
\newcommand{\thesisShortEn}{AI Tutor for \emph{Muallim al-Sani} and Model Adaptation}
\newcommand{\thesisAuthorRu}{Исмаилов Абдирахим Ерланович}
\newcommand{\thesisAuthorEn}{Abdirakhim Ismailov}               % латиница — [утвердить] как в паспорте
\newcommand{\thesisAuthorEnFull}{Ismailov Abdirakhim Erlanovich}
\newcommand{\thesisSupervisorRu}{Аскарбекулы Нурсултан}
\newcommand{\thesisSupervisorEn}{Nursultan Askarbekuly}
\newcommand{\thesisUniversityRu}{Университет Иннополис}
\newcommand{\thesisUniversityEn}{Innopolis University}
\newcommand{\thesisCity}{Иннополис}
\newcommand{\thesisCityEn}{Innopolis}
\newcommand{\thesisYear}{2027}

\hypersetup{pdftitle={Teaching Machines to Listen Like a Primer Teacher},pdfauthor={\thesisAuthorEn}}
\fancyhead[R]{\small\sffamily\color{inkgray}Research review · \thesisAuthorEn}
\setcounter{secnumdepth}{2}

\begin{document}
\thispagestyle{plain}
\begin{center}
{\sffamily\small\color{inkgray}\thesisUniversityEn{} · MSc in \thesisSpecialtyEn{} · Research review for the master's thesis}\\[10pt]
{\sffamily\LARGE\bfseries\color{accent}Teaching Machines to Listen Like a Primer Teacher\par}
\vspace{4pt}
{\sffamily\large\color{accent}Automatic pronunciation feedback for beginners reading Qur'anic Arabic with \emph{Muallim al-Sani}: state of the art, gaps and design implications\par}
\vspace{10pt}
{\color{gold}\rule{0.35\linewidth}{1pt}}\\[8pt]
{\large \thesisAuthorEn}\\[2pt]
{\small Supervisor: \thesisSupervisorEn}\\[2pt]
{\small\color{inkgray} Version 1.0 · October 2026}
\end{center}

\begin{keybox}[Abstract]
This review brings together six bodies of knowledge that a reading tutor for the primer
\emph{Muallim al-Sani} has to respect: the history and method of the book, the phonetic difficulties of
Russian-speaking beginners, research on computer-assisted pronunciation training, the state of automatic
assessment of Qur'anic recitation, ways to adapt a speech model to a new population with little data, and
the emerging practice of building software with AI coding agents.
It also summarises the legal and ethical constraints on collecting voice data. The evidence supports a
narrow claim: automatic feedback can multiply practice and first-line correction for beginners, but it
cannot yet replace a teacher. The best open model, Quran Muaalem, was trained on expert male reciters, and
no study has measured it on primer-level reading by second-language beginners or tried to adapt it to
them, although the methods --- low-rank fine-tuning, calibration, abstention with guarantees --- exist and
work with little data, provided that the evaluation is not spoiled by leakage. Explicit,
segment-level feedback is the condition under which pronunciation training works best, and it is what a
primer requires. Development with AI agents is productive but prone to gaming its own tests, which argues
for mechanical guard-rails rather than written rules. The review ends with a map of gaps and the research
questions they generate for the thesis.

\smallskip
\textbf{Keywords:} computer-assisted pronunciation training; mispronunciation detection and diagnosis;
Qur'anic recitation; Arabic as a second language; model adaptation; LoRA; calibration; conformal
prediction; Telegram bot; design science; AI coding agents; harness engineering.
\end{keybox}

\section{Introduction}

A child in a Tatar school at the end of the nineteenth century learned to read with a primer, and a
teacher who heard every syllable. Ahmad Hadi Maksudi's \emph{Muallim al-Sani} (``The Second Teacher'')
belongs to that tradition and is still in use: it is reprinted in Tatarstan and distributed by the
Muslim board of Kazakhstan \cite{SRC-020,SRC-040}. The adults who open it today are often
Russian-speaking, start from zero, and practise alone between lessons. What they lack is the listening
half of the method.

This review asks whether that half can be supplied by a machine, and under what conditions. It is the
research background for a master's thesis that builds a Telegram tutor following the book page by page,
recognises the learner's phonemes with an open model, explains errors in Russian, and adapts the model
to Russian-speaking beginners using the consented recordings it collects. Five questions guide the
review:

\begin{enumerate}[label=\textbf{Q\arabic*.}]
  \item What do the book and its teaching tradition require from a tutor?
  \item What makes Arabic reading hard for Russian- and Kazakh-speaking beginners?
  \item What does research say about automatic pronunciation feedback and how it should be designed?
  \item How good are current models for Qur'anic pronunciation assessment, and where are they untested?
  \item How can such a model be adapted to a new population when data are scarce, and how can the gain
    be measured honestly?
\end{enumerate}

A sixth, methodological question concerns the way the tutor itself is built: what is known about
developing research software with AI coding agents, and how to keep such development honest.

\paragraph{Method of the review} The review is narrative and scoping rather than systematic. Searches
were run in October 2026 over Google Scholar, arXiv, the ACL Anthology, the ISCA archive, Crossref,
publisher sites, eLibrary-indexed Russian journals, application stores and official regulator pages.
Every source received a stable key and a reliability grade: A for peer-reviewed or primary sources, B
for preprints, official documentation and authors' code, C for secondary sources such as blogs and
shop pages, used only for orientation. Bibliographic metadata of the key academic sources was checked
against Crossref or arXiv. The full registry, with what was taken from each source and where it is used
in the code, is kept in the project repository \cite{SRC-038}.

\section{The primer and its method}

\subsection{Two books by Maksudi}

Maksudi (1868--1941) wrote two primers that are often confused online. \emph{Mögallim äwwäl} (``The First
Teacher'', 1892) taught Tatar literacy in Arabic script; its sound-syllabic method and about thirty
editions before 1918 made it one of the most printed Tatar books of its time \cite{SRC-015,SRC-019}.
\emph{Muallim al-Sani} continued it: from knowing the letters to reading the Qur'an correctly, with the
basics of \emph{tajwid}. It is usually dated to the mid-1890s, but the year of the first edition has not
been confirmed by a primary bibliographic source \cite{SRC-018}. Modern texts descend from the Kazan
edition of 1917 \cite{SRC-024}. Because editions differ in length and content, the thesis fixes one
concrete edition, the 2010 Almetyevsk edition published under the Muslim board of Tatarstan
\cite{SRC-040}, and ties every exercise to its page and line.

\subsection{From letter names to sounds}

Tatar schools knew three ways of teaching literacy \cite{SRC-017}. The old letter-naming method
(\emph{usul-i harfiya}) made pupils memorise letter names, \emph{alif, ba, ta}, before assembling
syllables. The syllabic method (\emph{usul-i maddiya}) drilled syllables directly. The sound method
(\emph{usul-i sautiya}) of the Jadid reformers taught each letter as a sound and built syllables and
words consciously. The switch to the sound method is credited with reducing the time needed to learn to
read ``from several years to one school year'' \cite{SRC-016}. Damirov adds a detail that matters for an
Islamic audience: the reformer Z.~Rasulev justified the sound method by the tradition of \emph{tajwid}
itself rather than by European pedagogy \cite{SRC-016}.

Maksudi's innovation, as summarised by Fakhrutdinova and Ibragimov, lay in three choices: abandoning the
memorisation of letter names, attending to articulation and sound, and progressing gradually and
meaningfully \cite{SRC-015}. His fifteen didactic principles, published in 1906, fall into four groups.
Table~\ref{tab:principles} maps them onto design decisions for a tutor.

\begin{table}[ht]
\centering\small
\caption{Maksudi's didactic principles \cite{SRC-015} and their implications for a digital tutor.}
\label{tab:principles}
\begin{tabularx}{\linewidth}{@{}P{4.2cm}L@{}}
\toprule
Principle & Implication for the tutor \\
\midrule
Gradualness and continuity & Lessons in the book's order; a line counts as learned only when read
correctly; small steps (no lesson longer than a dozen lines). \\
Individualisation and activity & Learner's own pace and history; repetition of weak lines; any number
of retries. \\
Visual support and apperception & Reference audio before reading; the lesson's letter highlighted as in
the book; hints that relate new sounds to familiar Russian sounds. \\
Practice and continuous control & Every attempt checked immediately; each error explained; retry
straight away. \\
Humane treatment, no punishment & No penalties or rankings; a streak of days as encouragement, not a
threat. \\
Sound, not letter names & A card teaches the sound with a vowel sign; the name is reference only. \\
\bottomrule
\end{tabularx}
\end{table}

\subsection{Order of letters and the lesson cycle}

Two independent sources and the 2010 edition agree that the book does not follow the alphabet. It
introduces letters in groups, starting with the hamza and \ar{ر ز م}, letters with simple, mostly
non-joining shapes from which syllables such as \ar{زَرَ} and \ar{مَرَ} can be formed at once
\cite{SRC-021,SRC-040}. The order then continues
\ar{ت ن}, \ar{ي ب}, \ar{ك ل} and so on to \ar{ظ}, followed by long vowels, \emph{shadda} and
\emph{tanwin}, the article and \emph{sukun}, rules of \emph{mim} and \emph{nun}, peculiarities of Qur'anic
spelling, pausing, and finally short suras. The same order appears in an Uzbek edition and in a modern
course by the thesis supervisor \cite{SRC-070}, which suggests that it is Maksudi's own.

Every course found follows the same cycle, a digital form of \emph{talaqqi}: hear the model, repeat it
(the video course repeats each line three times), read alone, receive correction, and move on only after
an error-free reading \cite{SRC-021,SRC-022,SRC-058}. Two observations from practitioners shape the
design. Teachers advise against learning through transliteration, because ``sound cannot be conveyed on
paper'' \cite{SRC-021,SRC-022}. And the Kazakh edition states that the book does not replace a living
teacher. A tutor should therefore put the reference sound first, keep transliteration as a hidden hint,
and present itself as a practice partner, not as a substitute.

\section{The learners}

The intended learners are Russian-speaking adults in Russia and Kazakhstan who do not know the Arabic
letters or the sounds absent from Russian. No peer-reviewed study of their reading errors was found.
Table~\ref{tab:difficult} combines methodological sources for Russian speakers with studies of other
first languages \cite{SRC-023,SRC-034,SRC-035}. Pharyngeal, emphatic and interdental consonants are
difficult across first languages; vowel length and gemination are hard because Russian does not use
them to distinguish meaning.

\begin{table}[ht]
\centering\small
\caption{Typical difficulties of Russian-speaking beginners (methodological sources and other L1 studies).}
\label{tab:difficult}
\begin{tabularx}{\linewidth}{@{}P{3.2cm}L@{}}
\toprule
Sounds & What learners tend to do \\
\midrule
\ar{ح}, \ar{ع} & \ar{ح} becomes Russian [x] or merges with \ar{خ}, \ar{ه}; \ar{ع} becomes a glottal stop or
  disappears into the vowel. \\
\ar{ق}, \ar{غ}, \ar{خ} & \ar{ق} is fronted to [k]; \ar{غ} and \ar{خ} are confused. \\
\ar{ص}, \ar{ض}, \ar{ط}, \ar{ظ} & Emphatic consonants lose their emphasis and become [s d t z]. \\
\ar{ث}, \ar{ذ}, \ar{ظ} & Interdentals become [s]/[t], [z]/[d]. \\
Long vowels, \emph{shadda} & Long vowels are shortened; doubled consonants are not doubled. \\
\ar{ر} & Softened before front vowels. \\
\bottomrule
\end{tabularx}
\end{table}

The primer anticipates these confusions. Its ``pay attention'' lines contrast \ar{ق}/\ar{ك},
\ar{ث}/\ar{س}, \ar{ص}/\ar{س}, \ar{ط}/\ar{ت}, \ar{ح}/\ar{خ}, \ar{ع}/\ar{غ}, \ar{ض}/\ar{د} and
\ar{ذ}/\ar{ز}, which are minimal pairs in modern terms \cite{SRC-040}. For Kazakh speakers one testable
hypothesis follows from phonology: Kazakh has uvular [q] and [\ipa{ʁ}], so \ar{ق} and \ar{غ} may be easier for
them than for Russian speakers, while the other difficulties are shared. This remains a hypothesis until
data are collected.

\section{Computer-assisted pronunciation training}

\subsection{Does automatic feedback help?}

The strongest evidence comes from meta-analyses in second-language teaching. Ngo, Chen and Lai pooled
studies of automatic speech recognition in English pronunciation teaching and found a medium overall
effect, $g = 0.69$; effects were larger for segmental features than for prosody, and larger with
explicit corrective feedback than with an implicit transcript \cite{SRC-025}. Pronunciation instruction in
general has a large effect \cite{SRC-105}, so automatic feedback is not superior to teaching; it is a way
to deliver more of it. A systematic review of thirty studies by Amrate and Tsai adds two cautions: speech
recognisers often misjudge accented speech and give false alarms on intelligible pronunciation, and
supervised use, with a teacher or partner who can overrule the system, is therefore recommended
\cite{SRC-026}. Typical studies are quasi-experimental, with a few dozen participants over several weeks,
which sets a realistic bar for a master's pilot.

Two neighbouring lines of research support features beyond the core feedback loop. High-variability
phonetic training, in which learners hear contrasts from many voices, has a medium-to-large and lasting
effect on perception \cite{SRC-044}, with a smaller transfer to production \cite{SRC-045}. Distributed
practice is among the most robust findings in learning research \cite{SRC-107}, which justifies spaced
review of weak lines and a delayed post-test.

\subsection{How should feedback look?}

Classroom research on corrective feedback found that simply repeating the correct form (a recast) rarely
leads learners to correct themselves, while prompts that point to the error and ask for another attempt
do \cite{SRC-053}. Work on computer-assisted pronunciation training recommends feedback on a few
important errors rather than on everything \cite{SRC-054}. Reviews of formative feedback ask for
specific, simple feedback in small portions, with detail available on request \cite{SRC-055}, and
identify ``what to do next'' as the most useful element, while praise of the person adds little
\cite{SRC-056}. In the Qur'anic tradition the Iqro method and \emph{talaqqi} make the same point from
practice: the teacher signals where the mistake is and the learner corrects it by reading again
\cite{SRC-057,SRC-058}. Commercial tools highlight mistakes directly in the text \cite{SRC-046}.

Taken together, these sources suggest a feedback unit of three parts: \emph{what} went wrong and where
(the syllable is marked in the line), \emph{how} to say it (an articulation hint in the learner's
language), and \emph{what next} (listen, say the syllable slowly, read the line again). They also
suggest showing at most a few errors at a time and putting clear errors, such as a wrong letter, before
subtle ones such as vowel length \cite{SRC-052}.

\subsection{Telegram as a medium}

Telegram is the everyday messenger of the target audience and supports voice messages without an app
install. A four-week study found that a Telegram channel improved English pronunciation of Iranian
learners \cite{SRC-092}; no controlled study of a Telegram bot with automatic pronunciation checking was
found. The platform imposes its own rules: data used for machine learning may be collected only with the
user's explicit, active and revocable consent, and a privacy policy is required \cite{SRC-033}.

\section{Automatic assessment of Qur'anic recitation}

\subsection{The task and its metrics}

Mispronunciation detection and diagnosis compares what a learner said with what the text requires, at
the level of phonemes. The standard evaluation, adopted by L2-ARCTIC and by the Qur'anic benchmarks,
labels each reference phoneme as a true acceptance, true rejection (a real error that the system
found), false acceptance (a missed error) or false rejection (a false alarm). Precision is
TR/(TR+FR), recall is TR/(TR+FA), and F1 is their harmonic mean; diagnosis accuracy asks whether the
system also named the right substitute \cite{SRC-063,SRC-064}. For a beginner the two kinds of
mistake are not equal. A missed error delays learning, but a false alarm on a correct reading erodes
trust and motivation.

\subsection{Benchmarks}

QuranMB, introduced at Interspeech 2025, was the first open benchmark for Arabic mispronunciation
detection built on Qur'anic material \cite{SRC-004}. Its baseline systems had high recall and very low
precision. In the Iqra'Eval shared task the participating systems reached F1 around 0.40--0.44
\cite{SRC-043}. In the IQRA 2026 challenge on QuranMB.v2 the best system reached a phoneme-level F1 of
0.72 against a baseline of 0.44 \cite{SRC-031,SRC-086}. The winning systems are phoneme recognisers built
on self-supervised speech encoders with connectionist temporal classification, the same family as
Muaalem. The organisers state plainly that the data contain no second-language learners and call for
such data \cite{SRC-031}.

\subsection{Quran Muaalem}

Quran Muaalem is the most complete open system for diagnosing Qur'anic pronunciation errors
\cite{SRC-002,SRC-003}. It fine-tunes a Wav2Vec2-BERT encoder with eleven CTC heads: one for phonemes in a
dedicated Quran Phonetic Script and ten for articulatory attributes (\emph{sifat}) such as voicing,
emphasis and \emph{qalqala}. The training corpus has 848 hours and about 286 thousand utterances by
22 expert reciters, labelled largely automatically. Table~\ref{tab:muaalem} gives the reported results.

\begin{table}[ht]
\centering\small
\caption{Reported performance of Quran Muaalem \cite{SRC-002}.}
\label{tab:muaalem}
\begin{tabularx}{\linewidth}{@{}Lr@{}}
\toprule
Measure & Value \\
\midrule
Phoneme error rate, authors' test set (expert recitations) & 0.54\,\% \\
Average error rate across all eleven levels, same set & 0.21\,\% \\
Phoneme error rate, qdat\_bench (159 recordings with real errors) & 5.78\,\% \\
Average tajwid F1 on qdat\_bench (three rules) & 0.758 \\
F1 for \emph{ikhfa} on qdat\_bench & 0.453 \\
\bottomrule
\end{tabularx}
\end{table}

The ten-fold gap between expert test data and recordings with real errors is the most important number
for a beginners' tutor. The authors list the limitations themselves: training data come only from male
expert reciters in the Hafs reading, and qdat\_bench was annotated by one person on one verse
\cite{SRC-002}. The phonetiser expects Qur'anic text, so isolated letters and syllables of a primer need
canonical phoneme strings, and slow syllable-by-syllable reading is outside the training distribution.
The package also exposes per-phoneme confidence and attribute probabilities \cite{SRC-062}, which open
the way to a ``not sure'' outcome instead of a false verdict.

\subsection{Data on learners}

Data on non-Arab beginners are scarce. The closest public resource, a crowdsourced dataset of about seven
thousand recitations by non-Arab speakers from more than eleven countries, was created at Innopolis
University \cite{SRC-032}; it contains no Russian or Kazakh speakers. Work by the same group on
memorisation checking in a Telegram bot \cite{SRC-067} and on a taxonomy of recitation events
\cite{SRC-068} is the nearest neighbour of this thesis, with an important difference. Those studies
check \emph{what} was recited: which verses, in which order, with omissions and insertions, from
transcripts without vowel signs; vowels and \emph{tajwid} are explicitly out of scope. A primer tutor must
check \emph{how} a syllable was pronounced. The taxonomy is nevertheless useful because it separates
uncorrected mistakes from self-corrections and benign events, a distinction that beginners' recordings
will need.

\subsection{Applications}

Tarteel, the market leader, focuses on memorisation and flags word-level mistakes; it highlights them in
the text, a practice worth borrowing \cite{SRC-046}. A new wave of applications advertises AI checking
at the level of the primer, but their claims come from store descriptions and have not been evaluated
independently. Digital versions of \emph{Muallim al-Sani} exist as an app and an open web platform
\cite{SRC-005,SRC-006}, and the supervisor's course follows the same letter order while training
listening rather than speaking \cite{SRC-070}. No Russian-, Kazakh- or Uzbek-language tool that checks
primer-level reading by voice was found.

\section{Adapting a model to a new population}

\subsection{Fine-tuning on learner speech}

Mispronunciation detectors built on self-supervised speech encoders are usually obtained by fine-tuning
with CTC on annotated learner speech. Xu and colleagues showed on L2-ARCTIC that a fine-tuned wav2vec~2.0
outperforms earlier systems and keeps its quality when little labelled data is available
\cite{SRC-116}; Muaalem is built the same way, but was fine-tuned on experts \cite{SRC-002}. With a few
hours of speech, full fine-tuning of a 600-million-parameter model is risky: it overfits and may damage
what the model already knows. LoRA freezes the original weights and trains small low-rank matrices
added to the attention layers; the number of trainable parameters is orders of magnitude smaller, and
the original model is recovered by simply disabling the additions \cite{SRC-111}. Fine-tuning on a new
task can erase old skills --- catastrophic forgetting \cite{SRC-118} --- so the adapted model must also be
checked on its original domain, here expert recitation.

\subsection{Where the errors come from}

Beginners make more errors than experts, but there are still few of them for training, and particular
sound pairs are rare. The MDD literature augments data with synthetic errors by substituting phonemes in
the reference transcription as if the learner had erred \cite{SRC-119}; a standard speech-recognition
technique is speed perturbation \cite{SRC-120}. The primer offers a more natural source: its minimal
pairs (\ar{ق}/\ar{ك}, \ar{ص}/\ar{س}, \ar{ح}/\ar{خ} and others) mean that a correct reading of one card is
a real recording of a real error for the paired card. Such augmentation uses human audio rather than
synthesis, which the model recognises poorly (the project's preliminary check).

\subsection{Confidence, calibration and ``not sure''}

The second route is not to change the model but to learn where not to trust it. The classical Goodness
of Pronunciation score turns phoneme posteriors into a pronunciation score \cite{SRC-115}; a small
verifier can be trained on such features with a few hundred labelled examples. Neural networks are
overconfident, however, and their probabilities need calibration; the standard measure is the expected
calibration error \cite{SRC-112}. Selective classification lets a system abstain and describes the
trade-off with a risk--coverage curve \cite{SRC-117}. Conformal prediction provides distribution-free
guarantees under exchangeability \cite{SRC-113}, and conformal risk control chooses a threshold so that
the expected false-rejection rate stays below a chosen level \cite{SRC-114}. For the tutor this answers
the main danger that pronunciation-training reviews warn about: a false ``incorrect'' on a correct
reading.

\subsection{How not to fool ourselves}

Small data make inflated results easy. Tuning hyperparameters on the data on which quality is reported
biases the estimate upwards; nested cross-validation or a separate test set is required \cite{SRC-121}.
Kapoor and Narayanan found data leakage in hundreds of ML-based studies and proposed a taxonomy of it
\cite{SRC-122}; in speech, the main leak is the same speaker in training and evaluation. Hence the rules
of this work: a learner-disjoint split, a test set frozen before training, blind annotation of the test
set, and a single run on it. The supervisor's group formulated the same requirements for studies in
which models are tuned by AI agents \cite{SRC-067}. The dataset to be released comes with a datasheet
following Gebru and colleagues \cite{SRC-123}.

\section{Building research software with AI coding agents}

\subsection{From agent loops to harnesses}

A coding agent runs a loop of reasoning, acting with a tool and observing the result \cite{SRC-108}.
SWE-agent showed that the interface between agent and computer matters: a linter that rejects
syntactically invalid edits at the moment of editing substantially raised the share of solved tasks
\cite{SRC-109}. Practitioners now call everything around the loop a \emph{harness}: instructions,
hooks, sub-agents, tests, linters and logs. Böckeler distinguishes \emph{guides}, which steer the agent
before it acts, from \emph{sensors}, which observe and correct afterwards, and computational controls
such as type checkers from inferential ones such as an LLM reviewer \cite{SRC-029}. Anthropic's
engineering notes recommend external state, one feature per session, end-to-end checks and an explicit
prohibition on editing tests \cite{SRC-028}. No peer-reviewed definition of ``harness engineering''
exists yet; an exploratory study of 2\,853 repositories found that context files dominate, while
skills and sub-agents remain rare \cite{SRC-073}.

\subsection{What the evidence says}

Three findings are directly relevant. First, agents game evaluation. In ImpossibleBench, where the
specification contradicts the tests so that any pass means cheating, frontier models frequently passed
anyway, and for Claude models most cheating consisted of modifying the tests; making tests read-only
blocked this path \cite{SRC-027}. Reward hacking learned during training can even generalise into
sabotage \cite{SRC-082}. Second, written rules are weak. Repository context files such as
\texttt{AGENTS.md} did not improve, and sometimes reduced, task success while raising cost
\cite{SRC-074}; another study found savings in time and tokens with unchanged success \cite{SRC-075};
an ablation over 288 runs found no change in correctness \cite{SRC-077}; and a factorial study of 1\,650
sessions found that adherence to instructions declines within a session \cite{SRC-076}. Third,
mechanical feedback works. Feeding static-analysis findings back to the model reduced security issues
substantially over ten iterations \cite{SRC-078}, agents can write useful property-based tests
\cite{SRC-080}, and mutation-guided test generation at Meta produced tests that engineers accepted
\cite{SRC-079}.

Two cautions complete the picture. Experienced developers in a randomised study were 19\,\% slower with
AI assistance while believing they were faster \cite{SRC-030}, so productivity must be measured from
logs, not from impressions. And simulated students built from language models diverge systematically from
real ones \cite{SRC-083,SRC-085}, so an agent that ``plays the learner'' can find usability problems but
cannot replace a pilot.

\begin{keybox}[Implication]
Rules that matter should become checks that the environment enforces: frozen tests during
implementation, architecture contracts, type checks, protected human zones and a logged record of every
block. This is the principle the thesis applies to its own development and evaluates as RQ-4.
\end{keybox}

\section{Legal and ethical constraints}

A voice-based tutor about the Qur'an processes data that are personal, potentially biometric and
indicative of religious belief. In Russia the developer of a public bot is a data operator under Federal
Law 152-FZ; since September 2025 consent must be a separate document, and the database with citizens'
data must be located in Russia \cite{SRC-036}. Kazakhstan requires local storage and consent for
automated processing \cite{SRC-037}. Telegram's terms forbid collecting data for machine learning without
explicit, active and revocable consent \cite{SRC-033}. A defensible design therefore separates consent
to the service from consent to voice recording, never stores audio without consent, never builds voice
prints, uses pseudonyms in all research outputs, restricts research participation to adults and lets
learners delete everything. The tutor should use only the Arabic text of the mushaf and its own neutral
explanations; the Russian prose of the book is copyrighted and stays out of the product. Questions of
the data operator and ethics review are institutional and belong to the university.

\section{Synthesis: gaps and research questions}

Table~\ref{tab:gaps} sets out what is known, what is missing and how the thesis responds.

\begin{table}[ht]
\centering\small
\caption{Map of evidence and gaps.}
\label{tab:gaps}
\begin{tabularx}{\linewidth}{@{}P{2.6cm}LLP{1.4cm}@{}}
\toprule
Area & Known & Missing & Thesis \\
\midrule
Primer & Sound method, principles, letter order \cite{SRC-015,SRC-016,SRC-017} & Study of \emph{Muallim
al-Sani} as a Qur'an primer; any digital adaptation with checking & Artefact A \\
Learners & Difficult sounds from other L1 \cite{SRC-034,SRC-035} & Measured errors of Russian/Kazakh
speakers & RQ-1 \\
Feedback & Medium effect; explicit and segmental works best \cite{SRC-025,SRC-026} & Effect for Arabic
letters; trials of primer methods & RQ-3 \\
Models & Strong on experts; F1 0.72 on benchmark \cite{SRC-002,SRC-031} & Accuracy on beginners,
false-rejection rate, agreement with teachers & RQ-1 \\
Adaptation & LoRA, calibration, conformal guarantees \cite{SRC-111,SRC-112,SRC-114} & Adapting a
Qur'anic model to L2 beginners with honest evaluation on unseen learners & RQ-2 \\
Agents & Test gaming; rules decay; checks help \cite{SRC-027,SRC-076,SRC-078} & Whole-project evidence
on guarded harnesses & RQ-4 \\
\bottomrule
\end{tabularx}
\end{table}

From this map follow the four research questions of the thesis. RQ-1 asks how closely the
off-the-shelf model agrees with teachers on beginners reading the primer, and which reported errors
belong to the learner rather than the model. RQ-2, the central question, asks how much adapting the model
on the collected dataset --- a calibrated verifier with ``not sure'' and LoRA fine-tuning --- closes the gap
to a teacher on unseen learners without degrading expert recitation. RQ-3 asks whether explanations in
Russian improve the next attempt and retention compared with a bare correct/incorrect signal. RQ-4 asks
what a mechanically enforced harness prevents in AI-agent development and ML experimentation, and at
what human cost.

\section{Conclusion}

The literature leaves a clear and modest space for this work. The method of \emph{Muallim al-Sani} is
compatible with automatic feedback: it is sound-based, incremental and built on immediate correction.
The conditions under which pronunciation training works best, segmental targets and explicit feedback,
are exactly those of a primer. The model exists and is open. What is missing is measurement: how often the
model is right about a Russian-speaking adult reading a syllable for the first time, how often it rejects
a correct reading, whether its explanations help --- and how far it can be improved on the new
population's own data. A small, honestly annotated dataset, a frozen test set and a transparent table of
precision and recall against teachers before and after adaptation would be worth more than another claim
of ``AI checking''. The same dataset and test set become the starting point of doctoral work on
human-in-the-loop adaptation. In the same spirit, the development method should be judged by what its logs show, not
by how fast it felt.

\printbibliography[heading=bibintoc,title={References}]
\end{document}
